\documentclass[11pt]{article}
\usepackage[T1]{fontenc}
\usepackage{amsmath,amssymb,amsthm}
\usepackage[margin=1in]{geometry}
\usepackage[hidelinks]{hyperref}
\newtheorem{theorem}{Theorem}
\newtheorem{lemma}[theorem]{Lemma}
\newtheorem{corollary}[theorem]{Corollary}
\newcommand{\m}{\mathfrak m}
\newcommand{\Cc}{\mathcal C}
\newcommand{\Dd}{\mathcal D}
\title{Conductor power covers and failure of geometric envelopes\\
over arbitrary fields}
\author{ROOT working paper 425: scoped progress on P12.3}
\date{9 October 2026}
\begin{document}
\maketitle
\begin{abstract}
For every field and every exponent whose prime-to-characteristic part
exceeds one, an explicit irreducible nonnormal germ of dimension at
least three has a continuous signed primary-model obstacle without a
greatest continuous or decreasing-generated geometric minorant.
Its normalization has just one regular local centre. The construction
does not require roots of unity in the ground field: valuations with
finite residue extensions supply the distinct conductor lifts.
Every canonical degree has a uniform error at least $L-2$, for any
integer $L>2$, and failure persists under uniform perturbations.
The corresponding surface has continuous envelopes over every field
and for every exponent at the qualified normal-surface scope.
For a purely inseparable power, the full boundary map is instead
bijective and transports the higher-dimensional envelope question to
the regular normalization. These statements use the entire qualified
G331 geometric comparison and transport theorem. They do not solve
the general normal higher-dimensional envelope problem.
\end{abstract}

\section{The qualified geometric framework}
On the all-support normalized boundary $V_{\m}(R)$, let $\Cc_R$
be the uniform completion of negative effective finite real Cartier
evaluations on actual primary blowups, with nonnegative degree on
every contracted integral curve. Let $\Dd_R$ consist of actual
decreasing infima of $\Cc_R$, with the extended bottom convention.
The only incoming geometric correctness grant is the \emph{entire}
qualified G331 theorem, including every model, source, conductor,
radial, completion and historical-reading annotation.

We use its full ideal-profile comparison, finite-normalization
quotient map, positive descent, and exact primary-gauge transport.
In particular, inequalities for every finite real negative primary
profile extend to all of $\Cc_R$ and then to $\Dd_R$.
The former functions are continuous and nonpositive; the latter
are upper semicontinuous. Bijectivity of the full normalization
boundary yields an order/isometric identification of these cones
and equivalence of their continuous-envelope properties.
This is a substantive qualified geometric input, not a new definition
of a smaller cone. All subsequent ring and semivaluation calculations
are proved directly. The qualified excellent normal surface theorem
retained in the whole 422 reconstruction is used only in dimension two.

\section{The power-cover germs and their normalization}
Fix a field $k$, an integer $e\ge2$, and $d\ge3$. Put
$Z=(z_1,\ldots,z_{d-3})$ and define
\[
 A=k[u,y,s,Z]_{(u,y,s,Z)},\qquad t^e=u,
\]
\[
 S=A\oplus tA\oplus\cdots\oplus t^{e-1}A
   =k[t,y,s,Z]_{(t,y,s,Z)},
\]
\begin{equation}\label{eq:ring}
 R=A\oplus ytA\oplus\cdots\oplus yt^{e-1}A
   =k[t^e,y,yt,\ldots,yt^{e-1},s,Z]_{\m},
 \quad \m=(u,y,yt,\ldots,yt^{e-1},s,Z).
\end{equation}
The equalities include the indicated localizations, as follows.
Before localization the powers $1,t,\ldots,t^{e-1}$ give a free
module basis over $k[u,y,s,Z]$: grouping a polynomial by its
$t$-exponent modulo $e$ proves uniqueness and existence.
The displayed submodule $R$ is closed under multiplication, since
products with positive $t$-exponents have at least two factors $y$,
and reducing $t^e$ to $u$ preserves membership.
Both rings are finite integral algebras over $A$. Each has only one
maximal ideal over the maximal ideal of $A$: modulo it, $t$ in $S$
is nilpotent, while in $R$ all the $yt^i$ are nilpotent.
Their finite semilocalizations are therefore local, giving exactly
the rings in \eqref{eq:ring}.

The ring $R$ is a domain of dimension $d$, excellent as a local
finite-type $k$-algebra. The ring $S$ is regular and normal.
It is finite over $R$, and $t=(yt)/y$ belongs to the fraction field
of $R$. Thus the two fraction fields agree and $S$ is the finite
normalization of $R$, with a single marked local centre.

\begin{lemma}\label{lem:conductor}
The conductor is $\mathfrak c=yS\subset R$. Its quotients and gauge are
\[
 R/\mathfrak c=k[u,s,Z]_{(u,s,Z)},\quad
 S/\mathfrak c=k[t,s,Z]_{(t,s,Z)},\quad u=t^e,
 \qquad \m S=(t^e,y,s,Z).
\]
\end{lemma}
\begin{proof}
Certainly $yS\subset R$, and $(yS)S\subset R$.
Write an arbitrary element of $R$ as
$a_0+y\sum_{i=1}^{e-1}a_it^i$, $a_i\in A$.
If it belongs to the conductor, multiplying by $t$ still belongs
to $R$. The coefficient of $t$ in that product is $a_0$;
all coefficients of $t^i$, $i>0$, of an $R$ element are divisible
by $y$. Hence $a_0\in yA$, and the original element is in $yS$.
Conversely every such element is already in the conductor.
The quotient and gauge formulas now follow from the module descriptions.
Both reduced conductor quotients are domains.
\end{proof}

\section{Distinct lifts without ground-field roots of unity}
In characteristic zero assume $e\ge2$. In characteristic $p>0$,
write $e=p^r e_0$ with $p\nmid e_0$, and assume $e_0>1$.
There is a finite extension $\kappa/k$ containing an $e$th root
of unity $\zeta\ne1$: in positive characteristic the polynomial
$X^{e_0}-1$ has $e_0$ distinct roots in an algebraic closure,
and in characteristic zero use $X^e-1$. Adjoining one nontrivial
root gives a finite extension. No root in $k$ is required.

Define $\beta$ by the map $S\to k[[s]]$ sending
$t\mapsto s$, $y,Z\mapsto0$, followed by the $s$-adic order.
Define $\alpha$ similarly through $\kappa[[s]]$, sending
$t\mapsto\zeta s$, $y,Z\mapsto0$. These homomorphisms send
every local unit to a series with nonzero constant term, so define
centred real semivaluations on the stated local ring.
They are normalized by $\m S$, since $s$ has value one.
They are distinct: $\beta(t-s)=+\infty$ and $\alpha(t-s)=1$.
Nevertheless their restrictions to $R$ coincide. Both maps on
$R/\mathfrak c$ send $u\mapsto s^e$, $Z\mapsto0$, and on
$\mathfrak c$ have value infinity. Write $o$ for their common
target semivaluation. In particular, $o$ has nonminimal support
and the two lifts lie in the infinite-conductor stratum of the
\emph{same} regular normalization centre.

For real $r\ge1$, let $w_r$ be the semivaluation obtained by
$t\mapsto\zeta s$, $Z\mapsto0$ into $\kappa[y,s]$, with generic
monomial weights $w_r(y)=r$, $w_r(s)=1$. The same unit argument
defines it on $S$. Write $v_r=w_r|_R$.
Then $v_r(\m)=1$, $w_r(t-s)=1$, and
\begin{equation}\label{eq:convergence}
 v_r\longrightarrow o
 \quad\text{on every element of }R\text{ as }r\longrightarrow+\infty.
\end{equation}
Indeed, after substitution, the minimum of the weighted monomials of
a polynomial in $(y,s)$ tends to the order of its $y=0$ coefficient,
or to infinity if that coefficient vanishes. Unit denominators
have value zero. This proves convergence in the actual pointwise
boundary topology, not merely on the displayed generators.

\section{The obstacle and all-support positive competitors}
Fix an integer $L>2$. Let
\[
 J_1=(y(t-s))+\m^L,\qquad J_2=(y)+\m^L,
 \qquad F(v)=-v(\m)-v(J_1)+v(J_2).
\]
Here $y(t-s)=yt-ys\in R$. Both ideals are proper primary.
The function $F$ is an actual continuous signed primary-model
obstacle on a directed common blowup; its own positivity is not required.

Lift \emph{any} normalized target point to $w$ on $S$ using G331,
and put $Y=w(y)\ge1$, $a=w(t-s)\ge0$, allowing infinity. Then
\begin{equation}\label{eq:F}
 F(v)=-1-\{\min(Y+a,L)-\min(Y,L)\}\in[-L,-1].
\end{equation}
Only finite truncated quantities are subtracted. In particular
$F=-1$ on the entire infinite-conductor stratum.

For $n\ge1$ put
\[
 j_n=y(t-s)^n\in yS\subset R,\quad
 N_n=\m^{Ln}+(j_n),\quad T_n=\m^nN_n,
 \qquad q_n=-v(T_n)/n.
\]
All these are genuine proper primary ideals, and on every support
\begin{equation}\label{eq:q}
 q_n(v)=-1-\min(Y/n+a,L)\le F(v).
\end{equation}
For finite values the inequality follows from
\[
 0\le\min(Y+a,L)-\min(Y,L)\le\min(a,L)
 \le\min(Y/n+a,L).
\]
For infinite $Y$ or $a$ it follows directly from the finite truncations.

Before taking any limit, the blowup $\operatorname{Bl}_{T_n}R$
gives an effective Cartier divisor $E_n$ whose ideal is $T_n\mathcal O$.
Every standard domain chart embeds into the fraction field of $R$,
so its generating element is a nonzerodivisor.
The invertible sheaf $\mathcal O(-E_n)$ is relatively ample;
therefore $-E_n/n$ has strictly positive degree on every contracted
proper integral curve. Thus $q_n$ belongs to the actual geometric
core of $\Cc_R$, independently of any identification of its limit.
At the approaching points defined above,
\begin{equation}\label{eq:path}
 q_n(v_r)=-1-\min(r/n+1,L),\qquad
 \sup_{n\ge1}q_n(v_r)=-2.
\end{equation}

\section{A bound for the entire geometric cone}
Use the regular coordinates $(s,z,y,Z)$ on $S$, where $z=t-s$.
Let $\nu$ be the generic monomial valuation with weights
$\nu(s)=\nu(y)=\nu(z_i)=1$ and $\nu(z)=L$.
It is normalized by $\m S$: $\nu(t)=\nu(s+z)=1$, so
$\nu(t^e)=e$, while $\nu(y)=\nu(s)=1$.
Its restriction $\bar\nu$ satisfies
\begin{equation}\label{eq:deep}
 F(\bar\nu)=-L.
\end{equation}

For every proper primary ideal $I\subset R$ we have
\begin{equation}\label{eq:ideal}
 \bar\nu(I)\le o(I).
\end{equation}
The latter is finite because a power of $s$ lies in $I$.
Choose a finite generator attaining $e_I=o(I)$ and view it in
$S$ as a polynomial in $(s,z,y,Z)$ divided by a local unit.
Specializing $z=y=Z=0$ gives a nonzero polynomial in $s$ of
order $e_I$. Its numerator therefore has a nonzero pure monomial
$s^{e_I}$. The generic monomial valuation of that numerator is
at most $e_I$; distinct monomials do not cancel. Both denominators
have value zero. Unit profiles have value zero as well.
This proves \eqref{eq:ideal} for every profile entering the comparison.

Negating and summing, followed by the \emph{entire} qualified G331
real geometric comparison and actual decreasing-infimum extension, gives
\begin{equation}\label{eq:whole}
 g(o)\le g(\bar\nu)\qquad(g\in\Cc_R\cup\Dd_R).
\end{equation}
Every geometric minorant of $F$ consequently has $g(o)\le-L$.
The bound concerns the full geometric cone, not an ideal-only substitute.

\begin{theorem}\label{th:failure}
For the stated nonradicial exponents and every $d\ge3$, $F$ has
no greatest $\Cc_R$ or $\Dd_R$ minorant. The two raw suprema
agree, are finite, and fail upper semicontinuity at $o$.
Every continuous geometric minorant has uniform error at least $L-2$
from this common raw supremum $p$. In particular every canonical degree
has that same lower bound.
\end{theorem}
\begin{proof}
A greatest minorant would dominate every $q_n$, so by \eqref{eq:path}
its values along $v_r$ would be at least $-2$. Upper semicontinuity
and \eqref{eq:convergence} would give $g(o)\ge-2$, whereas
\eqref{eq:whole} and \eqref{eq:deep} give $g(o)\le-L<-2$.
The constant $-1-L$ is a core competitor, so both raw suprema
are finite. The same inequalities $p(v_r)\ge-2$, $p(o)\le-L$
prove failure of upper semicontinuity.

To identify the two raw suprema, let $g=\inf_i g_i\in\Dd_R$
be a decreasing presentation below the continuous obstacle $F$.
For each $\varepsilon>0$, the directed open sets
$\{g_i<F+\varepsilon\}$ cover the compact normalized boundary.
One index covers it, and $g_i-\varepsilon\in\Cc_R$ is a
minorant of $F$. Thus $g$ is below the raw $\Cc_R$ supremum
after letting $\varepsilon\downarrow0$. The reverse inequality
is immediate from $\Cc_R\subset\Dd_R$.

For a continuous minorant, $g(v_r)\to g(o)\le-L$ while
$p(v_r)\ge-2$, proving $\|p-g\|_\infty\ge L-2$.
The canonical ideals
\[
 A_m(F)=\{b\in R:v(b)+mF(v)\ge0\ \text{for every }v\},
 \qquad h_m=-v(A_m(F))/m
\]
contain $\m^{mL}$ and are proper. They are therefore primary.
Their defining inequalities give $h_m\le F$ everywhere, and
their actual blowups give the same negative effective relatively
ample Cartier presentations. Each $h_m$ is a continuous core
minorant, so the bound holds for \emph{every} $m\ge1$.
\end{proof}

If $\|F'-F\|_\infty<\varepsilon$ and $2\varepsilon<L-2$,
the actual competitors $q_n-\varepsilon\le F'$ force any greatest
minorant to be at least $-2-\varepsilon$ along $v_r$.
At $o$, \eqref{eq:whole} instead forces its value to be at most
$-L+\varepsilon$. Thus failure persists on this whole uniform open ball.

\section{The radicial case and the surface comparison}
\begin{theorem}\label{th:dichotomy}
In characteristic $p>0$, if $e=p^r\ge2$, the full normalization
boundary of the power-cover germ is bijective. Its geometric
continuous-envelope property is equivalent to that of $S$ under
the exact G331 order/isometric transport. No higher-dimensional
positive-characteristic envelope theorem is inferred.
For \emph{every} field and \emph{every} $e\ge2$, the corresponding
surface obtained by omitting $s$ and $Z$ has continuous envelopes
for every continuous signed obstacle at the qualified normal-surface scope.
\end{theorem}
\begin{proof}
Off the infinite-conductor stratum, $y$ has finite value and
$R_y=S_y$ because $t=(yt)/y$. The normalization lift is unique
there by the full G331 conductor description.
On the infinite-conductor stratum the ring map is
$k[u,s,Z]\subset k[t,s,Z]$, $u=t^e$.
If $e$ is a power of $p$, every source polynomial $f$ has
$f^e$ in the target polynomial ring, by the characteristic-$p$
power identity; coefficients need not have roots in $k$.
A target semivaluation has the unique extension
$w(f)=v(f^e)/e$. This formula preserves multiplication and
the ultrametric inequality, allows infinity, gives local units
value zero, and agrees on target elements since $v(a^e)=ev(a)$.
It extends to the localizations and respects the exact gauge
$(t^e,s,Z)$. Hence the conductor-boundary lift is unique and exists.
The full map is bijective, so the asserted equivalence follows
from qualified positive order/isometric transport.

For the surface, the conductor quotients are instead
$k[u]_{(u)}\subset k[t]_{(t)}$, with $u=t^e$.
The normalized conductor boundary downstairs has the sole valuation
$v(u)=1$: a nonzero element is a power of $u$ times a unit,
and a valuation killing $u$ cannot have finite normalized maximal
ideal. Upstairs the unique normalized valuation has $w(t)=1/e$.
Thus the full boundary is bijective for \emph{all} characteristics
and exponents, including separable ones. Its regular normal
normalization has dimension two, so the whole qualified
excellent-normal-surface envelope theorem applies. Transport gives
the claimed surface result.
\end{proof}

\begin{corollary}
Over every field there is an irreducible nonnormal surface with
universal continuous envelopes for which adding one affine coordinate
destroys this property at the corresponding origin.
One may take $e=2$ unless the characteristic is two, and $e=3$
in characteristic two. No ground-field roots of unity are necessary.
\end{corollary}
\begin{proof}
These choices have prime-to-characteristic part exceeding one.
Apply the surface assertion and the threefold failure theorem.
\end{proof}

\section{Completion, other gauges, and the original residual}
The same conclusions hold for the displayed completed rings.
Indeed replace $A$ by $k[[u,y,s,Z]]$ and $S$ by
$k[[t,y,s,Z]]$. Grouping power series by the exponent modulo $e$
gives the same finite free module descriptions, single regular
normalization centre, fraction field identity and conductor $yS$.
The series submodule defining $R$ is its maximal-ideal completion
because completion preserves the finite module description.
The generic monomial comparison still holds: a surviving pure
$s^{e_I}$ term bounds the weighted value, and all weights are positive.
For the path convergence, if the $y=0$ coefficient has finite order
$N$, every monomial with positive $y$-exponent has weight at least $r$,
so for $r>N$ it cannot affect that coefficient's leading value.
If the coefficient is zero, all values are at least $r$ and tend
to infinity. The maps into $\kappa[[y,s]]$ are continuous and
units retain nonzero constant terms. All-support formulas, actual
primary blowups, and the radicial extension identity therefore remain
valid without invoking a higher-dimensional resolution theorem.
For completed surfaces use the retained normal-surface theorem and
the identical DVR conductor-boundary argument.

Every displayed function is degree-one homogeneous on the unnormalized
cone. For another primary gauge $\mathfrak b$, use the positive
continuous $d(v)=v(\mathfrak b)$ and the exact G331 radial map.
Divide functions by $d(v)$; the obstruction gap at $o$ becomes
$(L-2)/d(o)$, and the canonical defining ideals are unchanged for
the transported obstacle. Uniform completion and decreasing
presentations transport with their full qualifications.

The failure is irreducible and \emph{nonnormal}. It gives no normal
higher-dimensional envelope counterexample. In radicial higher
dimension the problem is explicitly transferred to the regular normal
source at its actual characteristic hypotheses, rather than declared
solved. A universal criterion for arbitrary finite conductor covers,
general excellent normal dimension-at-least-three envelopes, and the
complete deeper recursive source audit remain open. Counts and the
original P12.3 statement remain unchanged. All primary literal bodies
and book context remain LOCALONLY; this is an unreviewed working paper.
\end{document}
